\chapter{Introduction}
\label{ch:intro}

\section{Background}

Dental panoramic radiography, also called orthopantomography (OPG), gives a single two-dimensional
view of the full dentition, the alveolar bone, the maxilla, the mandible and the surrounding
anatomy. Because one image covers most of the clinically relevant oral region, OPGs are used
routinely to screen for dental caries, periodontal bone loss, impacted teeth and other pathology.
They are also hard to read: local contrast is weak, structures overlap, the projection distorts,
and early lesions are subtle, so findings can be missed or reported inconsistently.

An OPG is unfamiliar to a computer-vision reader in two ways (Figure~\ref{fig:opgprimer}). It is
displayed as if the viewer faces the patient, so the patient's right side is on the image's left.
And teeth are named in the two-digit notation of the Fédération Dentaire Internationale (FDI): the
first digit is the quadrant (1 upper right, 2 upper left, 3 lower left, 4 lower right, in the
patient's frame) and the second the position counted from the midline (1 central incisor to 8 third
molar). Tooth 16 is therefore the upper-right first molar and 36 the lower-left one. Caries
dissolves mineral, so a lesion appears as a darker (radiolucent) region inside the bright tooth;
fillings and crowns are the brightest objects in the image.

\begin{figure}[H]
\centering
\includegraphics[width=0.86\textwidth]{figures/repo/R01_opg_primer.pdf}
\caption{Reading an OPG. One held-out DENTEX radiograph (\texttt{test\_120}) cropped to the teeth.
Dotted lines separate the quadrants Q1--Q4; the numbers are the FDI codes assigned by the
project's numbering stage. Yellow: the expert's seven caries annotations, on teeth 12, 16, 28, 35,
36, 37 and 46. Dashed: the eight-molar set \{16, 17, 26, 27, 36, 37, 46, 47\} named by the
image-blind constant of Chapter~\ref{ch:method}. Radiograph: DENTEX~\cite{hamamci2023dentex},
CC~BY-NC-SA~4.0.}
\label{fig:opgprimer}
\end{figure}

\section{Motivation}

The project builds an intelligent dental assisting system that moves from a full OPG to a
structured, tooth-wise report. The system is tooth-centric: instead of classifying the whole
radiograph as one unit, it decomposes the OPG into tooth-level regions and analyses each tooth with
task-specific models. The design follows clinical practice. Dentists reason tooth by tooth, and an
assistant is useful only if it can localise its findings and map them to a standard numbering
scheme. % {C070}

The survey carried out alongside the project reached a system-level conclusion that has shaped
every phase since: the dental-AI literature is strong on isolated sub-tasks and weak on complete
OPG-to-report workflows. A second observation came later, from the project's own measurements, and
it changed the direction of Phase 3. A caries system is usually scored by which teeth it marks as
carious, and that score can be earned partly without looking at the image at all, because caries is
concentrated on a few tooth positions. A fixed list of molars, the same for every patient, sets a
floor under any such score (Chapter~\ref{ch:problem}). Phase 3 therefore had two jobs: to finish the
end-to-end system, and to show what in its output a list that never reads the radiograph could not
produce.

\section{From Phase 2 to Phase 3}

Phase 2 (January--May 2026) delivered the disease-analysis modules separately: a two-stage ResNet18
caries cascade on DENTEX tooth crops, and a periodontal branch on the Tufts dental
database~\cite{panetta2022tufts} that included annotation parsing, tooth-level crop generation,
overlap-based labelling, a binary bone-loss classifier, a severity classifier and an exploratory
cemento-enamel junction (CEJ) module. Tooth detection, FDI numbering and report generation were
designed but not integrated, and no script took a full OPG to a report.

Phase 3 (August--December 2026) closed that gap and then tested it. Figure~\ref{fig:timeline}
summarises the progression, and Table~\ref{tab:phase-status} updates the Phase 2 implementation-status
table stage by stage.

\begin{figure}[H]
\centering
% Project progression, Phase 1 -> Phase 3. Components traced to: Phase 2 report (Ch. 1, Table 3.1);
% transfer.md section 1-2; this report Table 1.1.
\begin{tikzpicture}[node distance=6mm and 7mm,
  ph/.style={box,fill=#1,text width=44mm,minimum height=10mm,font=\sffamily\small\bfseries},
  items/.style={rectangle,rounded corners=2pt,draw=#1,line width=0.7pt,fill=#1!7,text width=44mm,
                font=\sffamily\scriptsize,align=left,inner sep=4pt}]
\node[ph=p2grey]   (p1) {Phase 1 \\[-1pt]{\normalfont\scriptsize Survey and design}};
\node[ph=p2indigo,right=of p1] (p2) {Phase 2 \\[-1pt]{\normalfont\scriptsize Jan--May 2026: modules}};
\node[ph=p2teal,right=of p2]   (p3) {Phase 3 \\[-1pt]{\normalfont\scriptsize Aug--Dec 2026: integrate, validate}};
\draw[arr] (p1) -- (p2); \draw[arr] (p2) -- (p3);
\node[items=p2grey,below=4mm of p1] (i1) {%
  \textbullet~Survey of OPG analysis\\
  \textbullet~Tooth-centric design\\
  \textbullet~Dataset roles chosen\\
  \textbullet~Pipeline-based taxonomy};
\node[items=p2indigo,below=4mm of p2] (i2) {%
  \textbullet~Caries cascade (ResNet18)\\
  \textbullet~Tufts parsing and crops\\
  \textbullet~Bone-loss binary + severity\\
  \textbullet~CEJ heuristic (exploratory)};
\node[items=p2teal,below=4mm of p3] (i3) {%
  \textbullet~Detector + FDI head + report\\
  \textbullet~Lesion detector (pre-registered)\\
  \textbullet~Image-blind constant theory\\
  \textbullet~Six cohorts, paired intervals\\
  \textbullet~CLI, Streamlit, Docker\\
  \textbullet~Phase 2 results corrected};
\draw[p2arrow,line width=0.6pt] (p1) -- (i1); \draw[p2arrow,line width=0.6pt] (p2) -- (i2);
\draw[p2arrow,line width=0.6pt] (p3) -- (i3);
\end{tikzpicture}

\caption{Project progression. Phase 1 surveyed the literature and fixed the tooth-centric design;
Phase 2 built the disease modules separately; Phase 3 integrated six stages into one pipeline,
validated it on six cohorts against an image-blind constant, and corrected two Phase 2 results.}
\label{fig:timeline}
\end{figure}

\begin{table}[H]
\centering
\caption{Implementation status at the end of Phase 2 and Phase 3. Phase 2 column: the status table
of the Phase 2 report. Phase 3 column: the repository on 2 November 2026.}
\label{tab:phase-status}
\small
\begin{tabularx}{\textwidth}{@{}L{4.6cm}L{3.3cm}Y@{}}
\toprule
Pipeline stage & End of Phase 2 & End of Phase 3\\
\midrule
Tufts data exploration and parsing & Complete & Complete; radiograph-level splits added\\
Tooth detection (Stage 1) & In progress, architecture only & Complete: YOLOv11l trained on three corpora, the default since 11 Oct 2026\\
FDI numbering (Stage 2) & Not yet integrated & Complete: ResNet18 head with box-position features and one-to-one Hungarian assignment\\
Caries detection (Stage 3a) & Complete (crop classifier) & Complete: a whole-image lesion detector, adopted under a pre-registered rule; the crop classifier is retained\\
Caries severity (Stage 3b) & Complete & Retained; not separately evaluated\\
Bone-loss / abnormality flag (Stage 4a) & Complete, macro-F1 0.9233 & Retrained with a radiograph-level split; the 0.9233 was a split artefact (Section~\ref{sec:res-stage4})\\
Bone-loss severity (Stage 4b) & Complete, macro-F1 0.7208 & Withheld from the report: retrained on mask overlap, it did not beat a tooth-position prior\\
CEJ-based estimation & Exploratory & Kept as a documented negative result\\
FDI-aware report generation & Not yet integrated & Complete: JSON report (schema \texttt{opg-report/1.6}) with an overlay image\\
End-to-end single-OPG inference & Not yet integrated & Complete: command-line tool, Streamlit demo and Docker image\\
\bottomrule
\end{tabularx}
\end{table}
% {C002} {C040} {C041}

\section{Novelties of Phase 3}
\label{sec:novelties}

The main contributions of Phase 3 are listed below; each is developed in the chapter named beside
it.

\begin{enumerate}[label=\arabic*.,leftmargin=1.6em]
\item \textbf{A complete six-stage pipeline.} One OPG in, one per-tooth JSON report out, through
  detection, numbering, caries localisation, severity, an abnormality flag and report fusion. It runs
  in 1.70~s per image on a laptop CPU with no GPU, as a command-line tool, a Streamlit demo and a
  Docker image (Chapters~\ref{ch:system} and~\ref{ch:impl}). % {C001} {C002}
\item \textbf{A theory of the tooth-level score.} The report is a chain of three decisions
  (detected, numbered, finding), and the caries score has an image-blind floor whose F1 has a
  closed form; the best constant contains exactly the codes whose prevalence exceeds half its F1
  (Chapter~\ref{ch:method}). % {C010}
\item \textbf{A constant-baseline evaluation protocol.} Two metric families, six cohorts (DENTEX, four
  Ningxia sites and CariesXrays), paired cluster-bootstrap intervals, random-effects pooling with a
  Hartung--Knapp--Sidik--Jonkman interval, and a comparison at matched output volume
  (Chapters~\ref{ch:method} and~\ref{ch:results}).
\item \textbf{A pre-registered change of task for caries.} Five repairs of the per-crop caries
  classifier failed; a whole-image lesion detector, adopted under a rule committed before training,
  led the constant on the pooled external sites while flagging fewer teeth (Section~\ref{sec:res-rq2}). % {C020}
\item \textbf{Outputs a constant cannot produce.} Tooth numbering with an error analysis, calibrated
  confidence and confidence-based abstention controlled by a shuffle test (Section~\ref{sec:res-rq3}). % {C034} {C035} {C036}
\item \textbf{Honest correction of Phase 2.} The bone-loss classifier was retrained with a
  radiograph-level split, the severity grade was withheld, and the shortcut signature of the
  per-crop caries classifier was measured (Section~\ref{sec:res-stage4}). % {C040}
\item \textbf{Traceable engineering.} Hash-pinned default checkpoints, 215 manuscript numbers asserted
  against saved artefacts, an end-to-end test suite and a reproducible container
  (Section~\ref{sec:impl-verify}). % {C050}
\end{enumerate}

\section{Objectives}

\begin{enumerate}[label=O\arabic*.,leftmargin=2.2em]
\item Integrate tooth detection, FDI numbering, caries analysis and an abnormality flag into one
  pipeline that produces a structured per-tooth report from a single OPG.
\item Run on commodity hardware without a GPU and package the system for demonstration.
\item Evaluate every stage on held-out and external data, beside an image-blind baseline, with
  intervals that respect the clustering of teeth within radiographs.
\item Change a model only under a rule written before the evidence, and report the rule's outcome
  whether or not it favours the change.
\item Correct any Phase 2 result that does not hold under a stricter split, in the open.
\end{enumerate}

\section{Scope}

The system is a research prototype. It is not a clinically validated device and makes no diagnosis;
every report says so (Section~\ref{sec:lld-report}). Its caries output is evaluated against expert
annotations of panoramic radiographs; its Stage-4 output flags proximity to an expert-marked
abnormality on the Tufts data and is not a measurement of periodontal bone loss. Training uses only
public datasets. One external cohort (Ningxia) is treated as locked test data and is never trained
or tuned on.

\section{Organisation of the Report}

Chapter~\ref{ch:problem} defines the problem and the research questions. Chapter~\ref{ch:lit} reviews
the literature. Chapter~\ref{ch:prs} specifies requirements. Chapters~\ref{ch:hld}
and~\ref{ch:lld} give the high- and low-level designs, and Chapter~\ref{ch:system} the detailed
design of every stage. Chapter~\ref{ch:method} sets out the datasets, the training recipes, the theory
of the image-blind constant and the evaluation protocol. Chapter~\ref{ch:impl} covers implementation
and pseudocode. Chapter~\ref{ch:results} reports and discusses the results, and
Chapter~\ref{ch:conclusion} concludes. Appendix~A lists definitions and acronyms; Appendix~B is the
user manual.
